\section{开源对齐的演进之路}

\subsection{指令微调模型的崛起（2023年春）}

2023年4月，第一批开源指令微调模型集中涌现：

\begin{figure}[H]
\centering
\includegraphics[width=0.95\textwidth]{slides-images/slide-022.jpg}
\caption{2023年4月的指令微调模型潮：Alpaca、Vicuna、Koala、Dolly 等\protect\footnotemark}
\end{figure}
\footnotetext{Slides 第23页。}

\begin{itemize}
    \item \textbf{Alpaca}：基于 LLaMA + 合成数据（Self-Instruct 方法）
    \item \textbf{Vicuna}：使用了 ShareGPT 数据——第一次让开放研究者获得了来自真实人类的对话数据
    \item \textbf{Koala}、\textbf{Dolly}：类似方法的变体
\end{itemize}

\begin{knowledgebox}{ShareGPT 数据的历史意义}
ShareGPT 是一个 Chrome 扩展，让用户可以分享 ChatGPT 对话。虽然存在法律灰色地带（用户数据被未经充分同意地记录），但这些数据对开源对齐研究至关重要——它是学术界首次大规模获取真实人类提示（prompt）的渠道。至今仍有模型将 ShareGPT 数据作为训练数据的子集。后来出现了更合规的替代品，如 LMSYS Chat 和 WildChat。
\end{knowledgebox}

\subsection{Open Assistant 与人类数据的稀缺}

\begin{figure}[H]
\centering
\includegraphics[width=0.95\textwidth]{slides-images/slide-024.jpg}
\caption{Open Assistant 项目：一次社区驱动的人类偏好数据收集尝试\protect\footnotemark}
\end{figure}
\footnotetext{Slides 第25页。}

Open Assistant 是一个 Discord 社区项目，由志愿者手动生成提示、回复和偏好对。Lambert 指出，虽然这个项目展示了人类数据的价值，但也暴露了其困难——参与者的反馈普遍是"再也不想做了"。遗憾的是，类似规模和质量控制水平的开源数据收集项目此后再未出现。

\subsection{第一个开源 RLHF 模型}

\begin{figure}[H]
\centering
\includegraphics[width=0.95\textwidth]{slides-images/slide-025.jpg}
\caption{StableVicuna：第一个开源 RLHF 模型，由 Carper AI 发布\protect\footnotemark}
\end{figure}
\footnotetext{Slides 第26页。}

2023年4月，Carper AI 发布了 StableVicuna——第一个使用标准 PPO 训练的开源 RLHF 模型。虽然它的性能优于 Vicuna，但\textbf{几乎没有人在此基础上继续构建}。Lambert 认为这揭示了一个重要教训：即使方法是开源的，如果没有易用的代码库和充足的数据，社区也不会跟进。

\subsection{Llama 2 安全过滤的反噬}

Llama 2 发布后出现了一个著名事件：当用户让模型"kill a Linux process"时，模型拒绝了——因为它将"kill"理解为暴力行为。这催生了一系列"uncensored"模型的开发潮流。

\begin{warningbox}{过度安全（Over-safety）的代价}
Llama 2 的过度安全问题并非有意的"审查"，而是 RLHF 训练数据的副作用。模型学到了过于保守的拒绝策略。这个案例说明：对齐不仅仅是"让模型更安全"，还需要在安全性和实用性之间找到平衡。过度拒绝同样是一种对齐失败。
\end{warningbox}

\subsection{Zephyr：DPO 的第一次大放异彩}

\begin{figure}[H]
\centering
\includegraphics[width=0.95\textwidth]{slides-images/slide-029.jpg}
\caption{Zephyr 模型：DPO 的首次重大成功\protect\footnotemark}
\end{figure}
\footnotetext{Slides 第30页。}

DPO 论文于 2023年5月发表，但直到2023年9月 Zephyr 模型的发布，社区才真正认可了 DPO 的实用性。Lambert 参与了 Zephyr 的开发，他回忆了两个关键成功因素：

\begin{enumerate}
    \item \textbf{新数据集}：UltraFeedback——由 OpenBMB 制作的合成偏好数据集，使用 GPT-4 进行标注
    \item \textbf{极低学习率}：$5 \times 10^{-7}$，比通常的 $3 \times 10^{-4}$ 低几个数量级
\end{enumerate}

\begin{warningbox}{DPO 需要极低的学习率}
Zephyr 的成功很大程度上依赖于一个反直觉的超参数选择：极低的学习率（$5 \times 10^{-7}$）。如果团队更早进行超参数搜索，DPO 可能提前几个月就被广泛采用。这提醒我们：即使算法本身是正确的，找到正确的超参数设置也至关重要——而这往往是通过大量试错才能实现的。
\end{warningbox}

\subsection{T\"ulu 2：DPO 规模化的验证}

\begin{figure}[H]
\centering
\includegraphics[width=0.95\textwidth]{slides-images/slide-030.jpg}
\caption{T\"ulu 2：第一个将 DPO 扩展到 70B 参数的模型\protect\footnotemark}
\end{figure}
\footnotetext{Slides 第31页。}

Lambert 转到 AI2 后，参与了 T\"ulu 2 项目。这个项目的核心贡献是：
\begin{itemize}
    \item 首次将 DPO 成功扩展到 \textbf{70B 参数}规模
    \item 使用相同的 UltraFeedback 数据 + 低学习率方案
    \item 在 MT-Bench 上达到 7.89 分（70B 版本）
    \item 正式引发了 DPO vs PPO 的学术讨论
\end{itemize}

从 Zephyr 到 T\"ulu 2 不到两个月。此后，DPO 模型如潮水般涌出——几乎每个发布开源模型的公司或团队都会推出一个"instruct"版本，通常基于 DPO 训练。

\subsection{本章小结}

从 2023年4月的第一批指令微调模型到 2023年底 DPO 的大规模应用，开源对齐经历了快速而曲折的发展。关键里程碑包括：ShareGPT 数据的出现、第一个 RLHF 模型、UltraFeedback 数据集的创建，以及 Zephyr 和 T\"ulu 2 对 DPO 的验证。整个过程中，\textbf{数据}始终是最稀缺也最关键的资源。

%% ========== Part 4: RewardBench ==========

